\subsection{Unknown tables and hidden keys}\label{sec:secret}
The following bounds quantify the cost of hiding the class. Their quantifier is one learner chosen independently of the flip table or key. That is different from choosing an architecture separately for each fixed class.

\begin{samepage}
\begin{proposition}[Unseen independent labels]\label{prop:unseen}
Fix the full-length template line $\ell_0:y=x+1$ and the uniform distribution on its $N$ points. Choose $A$ by independent fair incidence flips. Let a single learner, chosen independently of $A$, receive $T$ independent labeled samples from $(D,h_{\ell_0})$ and no further information about $A$. Then
\[
 \E_{A,\mathrm{sample},\mathrm{coins}}L_{D,h_{\ell_0}}(\widehat h)
 \ge\frac12(1-1/N)^T
 \ge\frac12(1-T/N).
\]
Thus a guarantee of average error at most $\epsilon$, or a guarantee for every flip table, requires $T\ge(1-2\epsilon)N$. If $\epsilon<1/4$ and $S\ge2$, such a uniformly successful $S$-parameter learner satisfies $TS\ge\dc(H_A)/3$ for every member of the flip family.
\end{proposition}
\end{samepage}
\begin{proof}
Condition on the sampled point indices, their observed labels, and the learner's random coins. The label at any unobserved point of $\ell_0$ remains an independent fair bit. The learner's prediction there has conditional error $1/2$. A fixed point is unobserved with probability $(1-1/N)^T$. Average over the $N$ points to get the first inequality. The second follows by induction from $(1-z)^T\ge1-Tz$ for $0\le z\le1$. Rearranging proves the sample requirement. For the final claim, $T>N/2$ and $S\ge2$ give $TS\ge N$, whereas $\dc(H_A)\le2N+1\le3N$ for every table.
\end{proof}

The lower bound is an average over tables, or a minimax statement against a learner that must work for all of them. It is false as a requirement for every individual fixed table. For example, the all-positive table needs no sample. The argument remains valid even if the learner is told $\ell_0$ and $D$: the unavailable information is the unobserved independent labels, not the line identity.

\begin{samepage}
\begin{proposition}[Hidden-key computational hardness]\label{prop:hiddenkey}
Under the PRF assumption of Theorem~\ref{thm:prf}, no probabilistic polynomial-time learner independent of $s$, with a polynomial-time evaluable output and only polynomially many labeled samples from the uniform distribution on $\ell_0$, has average error at most $1/2-\gamma(n)$ for uniform $s$, where $\gamma(n)\ge1/\operatorname{poly}(n)$. More precisely, a learner with $q(n)$ samples and advantage $\gamma(n)$ gives a PRF distinguisher of advantage at least
\[
 \gamma(n)-\frac{q(n)}{2N}.
\]
\end{proposition}
\end{samepage}
\begin{proof}
Given an oracle function, generate the learner's independent uniform point samples on $\ell_0$ and label them by its oracle values at the corresponding incidence encodings. Run the learner without giving it a key. Draw an independent uniform test point on $\ell_0$, query its label, and output one if the predictor agrees. In the PRF case acceptance is at least $1/2+\gamma(n)$. In the random-function case, conditional on the training transcript, a previously unqueried point has an independent fair label. A fresh test point repeats a training point with probability at most $q(n)/N$. Acceptance is therefore at most $1/2+q(n)/(2N)$. The simulation and one output evaluation are polynomial-time. Since $N=2^{(n-1)/3}$, the collision term is negligible, contradicting PRF security for inverse-polynomial $\gamma$.
\end{proof}

A guarantee for every key, or for all but a negligible fraction of keys, would imply the prohibited average guarantee. The proposition does not exclude an arbitrary algorithm specialized to one particular key. Nor does it automatically apply to infinite-precision real computation, which is not the finite-bit model in the cryptographic assumption.

\paragraph{The exact remaining SGD question.}
The original SGD premise permits the architecture, step size and horizon to be chosen for each class $H_A$. Consequently it does not itself impose independence from $A$, even though the weights are initialized from the prescribed Gaussian law. The two preceding propositions do not reverse those quantifiers. They prove obstructions for a \emph{single key-independent learner}, not an unconditional lower bound for every class-dependent benign-SGD design. Our SQ algorithms use the class through a table or an evaluable key; their generic differentiable simulations do not preserve the bias-free fully connected Gaussian-initialized logistic-SGD process. A negative answer for that process requires a compatible efficient representation and a proof of its dynamics for every marginal, not merely an oracle learner or a hidden-key lower bound. For an efficiently explicit class having a large rectangle ratio, Proposition~\ref{thm:classicalbarrier} also explains why the listed classical certificate parameters cannot establish superpolynomial sign-rank; it does not impose that barrier on every possible SQ-easy or SGD-easy class. Thus the original class-dependent benign-SGD question, and its stronger uniform key-independent variant, are stated separately rather than declaring the negative direction unreachable.
